\section{The Minimum Useful Founder Dossier}

\subsection{The Falsifiable Position Template}

A founder who has completed the audit should be able to state:

\begin{tcolorbox}[
  enhanced,
  colback=BayesPale!25!white,
  colframe=BayesBorder!60!white,
  boxrule=0.4pt,
  arc=2pt,
  left=6pt,right=6pt,top=4pt,bottom=4pt,
  before skip=3pt,after skip=4pt
]
\footnotesize\color{BayesInk}
I have demonstrable capability in \textbf{[3 task abilities]}, supported by \textbf{[artifacts/receipts]}. I understand \textbf{[2 or 3 workflows]} because \textbf{[direct experience]}. I can reach \textbf{[named personas with verified permission]}, and have \textbf{[cash, hours, tools]} for tests. My binding constraints are \textbf{[legal/debt limits]}. The opportunity I will test next is \textbf{[customer, pain, offer]}, because \textbf{[fit collision]}. The strongest objection is \textbf{[key unknown]}, which I will test by \textbf{[experiment by date]}.
\end{tcolorbox}

\textbf{The minimum useful founder-position dossier is more valuable than a pitch deck because it is falsifiable, constrained, and actionable today.}

\subsection{Worked Example: Retail Analytics Professional}

Consider a data analyst who spent 4 years optimizing operations for regional retail chains:

\begin{tcolorbox}[
  enhanced,
  colback=white,
  colframe=BayesBorder!80!white,
  boxrule=0.4pt,
  arc=2pt,
  left=5pt,right=5pt,top=4pt,bottom=4pt,
  before skip=3pt,after skip=4pt
]
\footnotesize\color{BayesInk}
\textit{I can clean messy transaction data, build replenishment dashboards, and explain stockout drivers to plant supervisors. I can show 2 deployed internal systems and client references. I understand small-chain replenishment routines, but do not yet understand their external vendor buying journey. I have access to 3 store owners willing to discuss pain, not 3 committed buyers. I can afford ₹25,000 in test spend, 12 hours weekly, and 1 weekend block. I cannot service weekday emergency support. My initial hypothesis is a fixed-scope stockout diagnostic for independent chains. My main uncertainty is whether store owners perceive reporting pain as worth paying an outside specialist to fix. I will interview 6 owners and offer 1 bounded paid diagnostic within 30 days.}
\end{tcolorbox}

\subsection{The Updating Founder Dashboard}

\begin{tcolorbox}[
  enhanced,
  colback=white,
  colframe=BayesBorder!80!white,
  boxrule=0.4pt,
  arc=2pt,
  left=3pt,right=3pt,top=3pt,bottom=3pt,
  before skip=3pt,after skip=5pt
]
\scriptsize
\setlength{\tabcolsep}{2pt}
\renewcommand{\arraystretch}{1.18}
\begin{tabularx}{\linewidth}{@{} >{\bfseries\color{BayesInk}}p{2.3cm} *{3}{c} >{\color{BayesSecondary}}X @{}}
\toprule
\textbf{Metric} & \textbf{D0} & \textbf{D30} & \textbf{D90} & \textbf{Diagnostic Purpose} \\
\midrule
Verified task skills & [ \ ] & [ \ ] & [ \ ] & Can we produce actual value? \\
Verified buyer talks & [ \ ] & [ \ ] & [ \ ] & Are we learning from reality? \\
Active decision buyers & [ \ ] & [ \ ] & [ \ ] & Can we test real demand? \\
Paid pilot deposits & [ \ ] & [ \ ] & [ \ ] & Did someone commit cash? \\
Hours per delivery & [ \ ] & [ \ ] & [ \ ] & Is delivery margin viable? \\
Cash spent vs loss cap & [ \ ] & [ \ ] & [ \ ] & Are we preserving survival? \\
Next falsification date & [ \ ] & [ \ ] & [ \ ] & What deserves focus now? \\
\bottomrule
\end{tabularx}
\end{tcolorbox}

\textbf{The dashboard prevents a founder from confusing motion with progress: track paid commitments, verified buyer conversations, and actual delivery hours rather than vanity attention.}
